\begin{abstract}
Las plataformas de observación terrestre como el Sistema Venezolano de Satélites Remotos (VRSS) proporcionan datos de alta resolución espacial pero con limitaciones en revisita temporal. Este artículo propone la integración de una red de nanosatélites (CubeSats) como sistema complementario para mejorar la resolución temporal y asegurar la continuidad de datos en aplicaciones de monitoreo ambiental, agrícola y de desastres. Se presenta una arquitectura híbrida que combina las capacidades de los satélites VRSS-1/VRSS-2 con constelaciones de nanosatélites en órbitas LEO optimizadas. Mediante modelado orbital, simulaciones de cobertura y algoritmos de fusión de datos multi-sensor, se optimiza la frecuencia de observación y la integridad temporal. Los resultados de simulaciones muestran mejoras significativas en la revisita (de días a horas) manteniendo resoluciones espaciales compatibles. Se discuten aspectos de interoperabilidad, procesamiento en borde y validación con datos reales. Esta aproximación ofrece una solución costo-efectiva y escalable alineada con estándares IEEE para sistemas satelitales distribuidos.
\end{abstract}